\documentclass[a4paper,12pt]{scrreprt}
    %% Used for changing geometry of the page
    %% Cover page text cannot overlay cover sketching/style
    %% https://ctan.org/pkg/geometry?lang=en
\usepackage{geometry}
    %% Changes language of some packages protocols
    %% e.g., when captioning images: Figure 1. -> Figura 1.
    %% https://ctan.org/pkg/babel?lang=en
\usepackage[portuguese]{babel}
    %% Used for special fonts
    %% Cannot be compiled with pdflatex
    %% https://ctan.org/pkg/fontspec?lang=en
\usepackage{fontspec}
    %% Arial FONT
    \setmainfont{Arial}

    %% More colors and color options
    %% https://ctan.org/pkg/xcolor?lang=en
    %% https://ctan.org/pkg/colortbl?lang=en
\usepackage{xcolor,colortbl}
    %% More tabular options, like dashed/dotted lines
    %% https://ctan.org/pkg/arydshln?lang=en
\usepackage{arydshln}
    %% List of acronyms
    %% https://ctan.org/pkg/nomencl?lang=en
\usepackage[intoc]{nomencl}
    %% Must be called to init nomencl environment
    \makenomenclature
    %% More images options/settings
    %% https://ctan.org/pkg/graphicx?lang=en
\usepackage{graphics}
    %% Defining subdirectories to image path enviornment
    %% \graphicspath{{sub1}{sub2}...{subN}}
    \graphicspath{{images}}

    %% used to handle cross-referencing commands in LaTeX to produce hypertext links in the document
    %% https://ctan.org/pkg/hyperref?lang=en
\usepackage{hyperref}
    %% math environments
    %% https://ctan.org/pkg/amsmath?lang=en

    %% settings
    \hypersetup{
        colorlinks,
        citecolor=black,
        filecolor=black,
        linkcolor=black,
        urlcolor=black
    }

\usepackage{amsmath}
    %% Defining backgrouns, used to make the cover
    %% https://ctan.org/pkg/background?lang=en
\usepackage[some]{background}
    %% Used to make drawings or complex graphics
    %% http://pgf.sourceforge.net/pgf_CVS.pdf
\usepackage{tikz}
    %% Tikz library to point operations ((x1,y1) + (x2,y2))
    \usetikzlibrary{calc}

%% Defining sfdefault font and default font for document
\renewcommand{\familydefault}{\sfdefault}

%==========================================================================
% DOCUMENT
%==========================================================================

\begin{document}

\pagenumbering{gobble}

%% Costume made cover
%% From there you can use \makecover command to build the cover
%% Blue cover color
\definecolor{titlepagecolor}{RGB}{54,95,145}

%==========================================================================
% COLORED BAR ON THE LEFT SIDE
%==========================================================================

\backgroundsetup{
    scale=1, 
    angle=0, 
    opacity=1,
    contents={
        \begin{tikzpicture}[remember picture,overlay]
            \path [fill=titlepagecolor] 
                (current page.north west) -- ($(current page.north west) + (5,0)$)
                -- ($(current page.south west) + (5,0)$)-- (current page.south west); 
            \node[color=white] at ($(current page.south west) + (3,4)$) {\bfseries {\fontsize{120}{60} \textsf{LI}}};
            \node[color=titlepagecolor] at ($(current page.south west) + (5.8,4)$) {\bfseries {\fontsize{120}{60} \textsf{4}}};
        \end{tikzpicture}
    }
}

%==========================================================================
% TITLE PAGE INFO
%==========================================================================

%% Changes values in this field to show information in the cover and back cover about your team/project


%% TITLE
\title{CDC – Consultoria de Detetives Christie}

%% AUTHORS
\author{
    Afonso Dionísio Santos (A10426) \\
  \quad
    Ana Filipa Cruz Pinto (A96862) \\
  \quad
    Carlos Humberto da Silva Ferreira (A89509) \\
  \quad
    Flávia Alexandra Silva Araújo (A96587) \\
  \quad
    Miguel Torres Carvalho (A95485)
}

%% Date

\date{\today}

%% Course
\newcommand{\Course}{Licenciatura em Engenharia Informática}

%% Department
\newcommand{\Department}{Escola de Engenharia}

%% UniName
\newcommand{\UniName}{Universidade do Minho}

%% UniPic
\newcommand{\UniPic}{\includegraphics[scale=0.09]{images/uminho.png}}

%% University 
\newcommand{\University}{
    \begin{flushleft}
        \UniPic
    \end{flushleft}
    \textcolor{gray}{\small\textbf{\textsf{\UniName}}}\par
    \textcolor{gray!80!white}{\small{\textsf{\Department}}}\par
    \textcolor{gray!70!white}{\small{\textsf{\Course}}}
}

%% UC
\newcommand{\UC}{
    \begin{flushleft}
        \par\textcolor{titlepagecolor}{  \LARGE\textbf{\textsf{Unidade Curricular de \\ Laboratórios de Informática IV}}}
    \end{flushleft}
}

%% School Year
\newcommand{\SchoolYear}{
    \small{\textsf{Ano Letivo de 2023/2024}}}


%% Define new command to show title, author and date
\makeatletter
\let\Title\@title
\let\Author\@author
\let\Date\@date
\makeatother

%==========================================================================
% CLASSIFICATION SECTION 
%==========================================================================

%% School Year
\newcommand{\ReceptionDate}{}
%% Responsible
\newcommand{\Responsible}{}
%% Evaluation
\newcommand{\Evaluation}{}
%% Observations
\newcommand{\Observations}{}





%% MAKETEMPLATE
\newcommand{\makecover}{

%==========================================================================
% BEGIN COVER PAGE 
%==========================================================================

%% Removes page number on footer
\thispagestyle{empty}

%% No indentation 
\setlength{\parindent}{0em}

%% Put Background defined on \backgroundsetup, in this page
\BgThispage

%% Changing geometry to prevent overlay with text
%% At the end of back cover, geometry is default with \restoregeometry
\newgeometry{top=5cm,left=6cm,right=3cm,bottom=2cm}

%% builds university info defined previously
\University
\vspace{1cm}
%% builds curricular unity info defined previously
\UC
%% builds school year info defined previously
\SchoolYear

\vspace*{5cm}
%% bigger space (i think its the default one) between paragraphs 
\setlength{\parskip}{1em}

%% builds title info defined previously
\par\textbf{\textsf{\huge\Title}}
\vspace{1cm}
%% builds author(s) info defined previously
\par\Author

\vspace{0.5cm}

%% builds date info defined previously
\par\Date
\restoregeometry
\pagebreak

%==========================================================================
% END COVER PAGE 
%==========================================================================

%==========================================================================
% BEGIN BACK COVER PAGE 
%==========================================================================

%% Removes page number on footer
\thispagestyle{empty}

% Changing look of lines in tabular environment 
% Dashed -> dotted 
%% length of dashes
\setlength\dashlinedash{0.3pt}
%% space between dashes
\setlength\dashlinegap{1.5pt}
%% width of dashes
\setlength\arrayrulewidth{1.1pt}


%% This values can be changed in the preamble
\begin{flushright}
\begin{tabular}{ :p{4cm}:p{4cm}: } 
\hdashline
Data de Receção & \ReceptionDate \\ [2ex]
\hdashline
Responsável & \Responsible \\ [2ex]
\hdashline
Avalição & \Evaluation \\ [2ex]
\hdashline
Observações & \Observations \\ [7ex]
\hdashline
\end{tabular}
\end{flushright}


\vspace{9cm}
\begin{flushleft}

%% builds title info defined previously
\par\textbf{\textsf{\huge\Title}}
\vspace{1cm}
%% builds author info defined previously
\par\hspace{0.25cm}\Author

\vspace{0.5cm}

%% builds date info defined previously
\par\Date
\end{flushleft}

\pagebreak
%==========================================================================
% END BACK COVER PAGE 
%==========================================================================
}


% builds the cover
\makecover

%% smaller footer and header size
\newgeometry{top=3cm,left=3cm,right=3cm,bottom=4cm}

%==========================================================================
% BEGIN OPCIONAL DEDICATÓRIA
%==========================================================================

% \clearpage
% \begin{center}
%     \thispagestyle{empty}
%     \vspace*{\fill}
%
%     $<<$/opcional Dedicatória$>>$
%
%     \vspace*{\fill}
% \end{center}
% \clearpage

%==========================================================================
% END OPCIONAL DEDICATÓRIA
%==========================================================================

%==========================================================================
% BEGIN ABSTRACT PAGE
%==========================================================================

%% Abstract name: \Large font size, flushed left and paragraph skip before abstract content
\renewenvironment{abstract}
 {\par\noindent\textbf{\Large\abstractname}\par\bigskip}
 {}

\begin{flushleft}
\begin{abstract}
    No âmbito da UC Base de Dados, lecionada pelo regente, Professor Orlando Belo, visamos a
    realização de um projeto que consiste na modelação, desenvolvimento e implementação de um
    Sistema de Base de Dados.
    \par Tendo em consideração o tema deste ano - uma agência de detetives -, decidimos explorar
    uma agência liderada por Agatha Christie, de nome CDC - Consultoria de Detetives Christie.
    Relativamente à modelação e desenvolvimento do nosso projeto, iremos dividi-lo em partes:
    começaremos de forma abstrata, transformando os factos
    \par \textbf{Área de Aplicação}: \textcolor{red}{
        <<Identificação da Área de trabalho. Por exemplo: Desenho e arquitectura de Sistemas de Bases de Dados.>>
    }
    \par \textbf{Palavras-Chave}: \textcolor{red}{
        <<Conjunto de palavras-chave que permitirão referenciar domínios de conhecimento, tecnologias, estratégias, etc., directa ou indirectamente referidos no relatório. Por exemplo: Bases de Dados Relacionais, Gestão de Índices, JAVA, Protocolos de Comunicação.>>
    }
\end{abstract}
\end{flushleft}

\pagebreak

%==========================================================================
% END ABSTRACT PAGE
%==========================================================================

%==========================================================================
% BEGIN INDEXES PAGES
%==========================================================================

%% Changes table of content name
%% Portuguese babel default : "Conteúdo"
%% Personally I prefer "índice"
\renewcommand{\contentsname}{Índice}
\renewcommand{\listfigurename}{Índice de Figuras}
\renewcommand{\listtablename}{Índice de Tabelas}

\tableofcontents

\pagebreak

\listoffigures

\pagebreak

\listoftables

\pagebreak

%==========================================================================
% END INDEXES PAGES
%==========================================================================

%==========================================================================
% BEGIN INTRODUCTION
%==========================================================================

%% Starting page numbering here
\pagenumbering{arabic}

\chapter{Definição do Sistema}
    \section{Contexto de Aplicação}
    Agatha Christie, uma figura proeminente no mundo dos detetives, criou a sua própria agência
    no final dos anos 90 após concluir que a sua carreira como detetive privada não ia ser
    suficiente para vingar-se do mundo sujo e curioso do crime.
    \par A sua agência começou como algo discreto - um escritório na periferia de Londres,
    constituído por Agatha - gerente e secretária, a cara da Consultoria de Detetives Christie
    (CDC) -, e mais três detetives, responsáveis por resolver os casos dos clientes que recorriam
    à agência nos seus momentos de aflição.
    \par Não obstante, nos últimos três anos, houve um crescimento exponencial de casos, visto que a
    sua agência tornou-se renomada devido a alta variedade de casos que é capaz de solucionar -
    desde os mais “mundanos”, como casos de infidelidade e perseguições, até aos mais
    “mórbidos”, como homicídios e desaparecimentos. E, visto que Agatha é fascinada pelo avanço
    tecnológico, a sua consultoria também é exemplo de vanguarda na solução de cibercrimes.
    \par Como tal, toda esta nova popularidade acrescida levou a que Agatha contratasse um novo
    estagiário, aumentando a sua equipa e procurando conseguir prepará-lo para a subida de
    casos que a agência enfrentava a todo o vapor.
    \section{Motivação e Objetivos do Trabalho}
    A CDC enfrenta um aumento significativo na demanda pelos seus serviços de investigação devido à
    sua reputação crescente e à diversificação dos casos que lida. Infelizmente, Agatha sentiu a sua
    valiosa agência a sofrer complicações a partir do momento em que decidiram aceitar um maior número
    de casos - O aumento na procura por serviços de investigação levou a uma sobrecarga nos sistemas de
    gerenciamento de casos existentes, e os registos físicos que ela mantinha desde o início da sua
    agência não lhe permitiam atribuir com rapidez suficiente os seus detetives aos casos, e muitas das
    informações cruciais, como pistas ou relatos de testemunhas, já haviam sido perdidos ou duplicados
    no passado, o que fazia Agatha temer que a sua agência acabasse por ficar com uma má reputação.
    \par Com a sua mente analítica e perspicaz, ela reconheceu que a chave para resolver este mistério
    organizacional estava  na modernização tecnológica, nomeadamente, na implementação de um Sistema de
    Base de Dados que possa lidar com a crescente quantidade de informações e casos de forma eficiente e
    escalável, bem como gerenciar e organizar as informações relacionadas aos casos, clientes, evidências
    e suspeitos. Este projeto visa atender a essa demanda e proporcionar à CDC as ferramentas necessárias
    para continuar a oferecer serviços de alta qualidade e eficácia, assim garantindo o sucesso da agência
    e aliviando as preocupações de Agatha sobre a popularidade acrescida.
    \par Por conseguinte, os objetivos principais que pretendemos alcançar com o desenvolvimento deste SBD são os seguintes:
    \begin{itemize}
        \item \textbf{Escalabilidade:} À medida que a CDC cresce e enfrenta um aumento contínuo na procura pelos
            seus serviços, é essencial ter um sistema que possa escalar para atender às necessidades em constante
            evolução da agência. Um Sistema de Base de Dados escalável pode crescer junto com a CDC, garantindo que
            ela permaneça ágil e adaptável às mudanças no mercado, sem falhas ou confusões no sistema.
        \item \textbf{Centralização dos Dados:} Com um Sistema de Base de Dados, os dados relacionados a
            casos, clientes, evidências e investigações podem ser acedidos rapidamente num local centralizado de
            forma rápida e eficiente, o que permite uma colaboração mais eficaz entre os detetives e facilita a
            tomada de decisões informadas.
        \item \textcolor{red}{
                \textbf{Eficiência Operacional:} Com o aumento do volume de casos, os métodos manuais de organização de informações tornaram-se cada vez mais ineficientes. Um Sistema de Base de Dados pode automatizar diversas tarefas, como armazenamento, recuperação e atualização de dados, libertando tempo e recursos dos funcionários para se concentrarem na própria investigação.
            }
        \item \textbf{Precisão e Consistência:} Os registos físicos estão sujeitos a erros humanos, como duplicação e perda de dados. Um Sistema de Base de Dados garante precisão e consistência nas informações, ajudando assim a evitar erros e inconsistências que possam comprometer a qualidade do trabalho da CDC.

        \item \textbf{Segurança de dados:} Os registos físicos podem ser facilmente acedidos por qualquer pessoa
            que os encontre. Isso inclui funcionários não autorizados ou intrusos, provocando falsificações, destruição
            acidental e intencional de provas. Com a implementação de Sistema de Bases de Dados, existe um maior
            controlo de acesso.
        \item \textbf{Controlo de Despesas:} Com a quantidade de informação que ocorre durante um caso, alguns registos de despesas podem ser esquecidas, por isso é importante saber qual foi o custo de um caso.
    \end{itemize}
    \section{Análise da Viabilidade do Processo}
        A viabilidade de um projeto de desenvolvimento de software depende da habilidade de compreender e satisfazer as
        demandas do mercado e dos utilizadores. Isto requer um planeamento cuidadoso e eficiente para garantir a
        entrega de um produto confiável e de alta qualidade. E, ao seguir uma abordagem metódica, o projeto pode
        maximizar as suas chances de sucesso ao atender às expectativas e necessidades do público-alvo de forma eficaz.
        \par Consideramos que o desenvolvimento de um SBD para a CDC é bastante viável, pois este garantirá uma série de benefícios para a agência, nomeadamente:
        \begin{itemize}
            \item \textbf{Melhor Gestão de Funcionários:} Com um Sistema de Base de Dados, existe uma maior facilidade para identificar que funcionários estão ocupados ou disponíveis, possibilitando uma alocação mais rápida dos mesmos aos casos, bem como uma melhor assistência destes conforme necessária.
            \item \textbf{Melhor Gestão da Consultoria:} Ao conhecer os custos de cada caso, é possível otimizar os recursos financeiros, planear orçamentos mais precisos e tomar decisões estratégicas fundamentadas para maximizar a eficiência e rentabilidade da empresa.
            \item \textbf{Melhorar a Qualidade de Serviço e de Bem Estar no Trabalho:} Um Sistema de Bases de Dados
                promoverá um melhor bem estar aos seus funcionários, evitando buscas intensivas ao sistema de
                informação já recolhida, consequentemente melhorando a qualidade do serviço, significativamente.
            \item \textbf{Resolver a Sobrecarga:} Devido aos dois pontos referidos anteriormente, os funcionários
                serão capazes de resolver um caso com mais eficiência e rapidez, ficando disponíveis mais rapidamente.
                Como tal, a sua produtividade vai aumentar e vai ficar a par da nova enchente de casos.
            \item \textbf{Segurança Acrescida:} Com um Sistema de Base de Dados, existe um maior controlo de acesso relativamente a informações cruciais aos casos, o que garante a inexistência de adulteração ou destruição de provas. Com isto, tem-se a certeza que as informações presentes nos registos são as originais e não foram acedidas por intrusos.
        \end{itemize}
        \par Considerando esses fatores, fica claro que o projeto de implementação do Sistema de Base de Dados é
        altamente viável e trará benefícios substanciais para a CDC, principalmente a níveis financeiros, de
        organização de dados e serviços, e, a longo prazo, de crescimento contínuo no mercado de detetives particulares.
    \section{Recursos e Equipa de Trabalho}
        \subsection{Recursos Humanos:}
            \begin{itemize}
                \item Funcionários da Consultoria (Detetives, estagiários e gerente);
                \item Clientes (Vítimas e os seus familiares, etc);
                \item Equipa de desenvolvimento.
            \end{itemize}
        \subsection{Recursos Físicos:}
            \begin{itemize}
                \item Computadores;
                \item Conexão à \textit{Internet};
                \item Servidor.
            \end{itemize}
        \subsection{Recursos Digitais:}
            \begin{itemize}
                \item Sistemas Operativos: \textit{Windows} 11 e \textit{Linux} (\textit{Ubuntu} 22.04.3 \textit{LTS});
                \item \textit{Google Drive};
                \item \textit{GitHub};
                \item \textit{LaTeX};
                \item brModelo (v3.31);
                \item \textcolor{red}{(...) vamos preenchendo isto à medida que vamos usando ferramentas}.
            \end{itemize}
        \subsection{Equipa de Trabalho:}
            \begin{itemize}
                \item \textbf{Perssoal Interno:}
                \begin{itemize}
                    \item \textbf{Agatha Christie:} Funcionamento da agência, atendimento a clientes,
                        validação de serviços, atribuição de casos aos agentes, depoimento de informações
                        cruciais ao projeto.
                    \item \textcolor{red}{\textbf{Detetives efetivos:}}
                    \item \textcolor{red}{\textbf{Detetives estagiários:}}
                \end{itemize}
            \item \textbf{Pessoal Externo:} \textcolor{red}{(para continuar ao longo do tempo)}
                \begin{itemize}
                    \item \textbf{Afonso Santos:}
                    \item \textbf{Ana Pinto:}
                    \item \textbf{Carlos Ferreira:} Levantamento de Requisitos
                    \item \textbf{Flávia Araújo:}
                    \item \textbf{Miguel Carvalho:} Levantamento de Requisitos
                \end{itemize}
            \end{itemize}

    \section{Plano de Execução do Projeto}
        \par Para assegurar uma implementação eficiente e eficaz do sistema de bases de dados, foram realizadas reuniões e conversas com a Sra. Christie, detetives e estagiários envolvidos no projeto. Com base nessas interações, foi estabelecido o seguinte cronograma de execução:
        \par
         \textcolor{red}{
            <<TODO: adicionar diagrama de Gantt>>
        }
        \par No cronograma referido, definimos o tempo e os intervenientes nas           várias fases do projeto.

    \section{Estrutura do Relatório}
        \textcolor{red}{
            <<Após a leitura da introdução de um relatório é "simpático" apresentar uma breve descrição daquilo que se vai encontrar nos demais capítulos do relatório.>>
        } \\
        \textcolor{red}{[isto faz-se no fim]}
%==========================================================================
% END INTRODUCTION
%==========================================================================

%==========================================================================
% BEGIN LEVANTAMENTO E ANÁLISE DE REQUISITOS
%==========================================================================

\chapter{Levantamento e Análise de Requisitos}
    \section{Método de Levantamento e de Análise de Requisitos Adotado}
            O processo de definição de requisitos começa com uma reunião com a equipa para selecionar diversas estratégias de levantamento, de modo a captar toda a informação necessária para sustentar adequadamente e logicamente uma melhor esquematização e composição de requisitos. Das várias estratégias discutidas, assumimos que utilizamos estas:
            \begin{itemize}
            \item \textbf{Reuniões} presenciais ou onlines com a Sra. Christie e os seus Detetives, tanto em grupo quanto individualmente com o propósito de identificar os diversos processos operacionais que ocorrem na consultoria e classificá-los.
            \item \textit{\textbf{Emails}}, às vezes as melhores ideias podem surgir subitamente, e como os funcionários têm uma vida bastante atarefada, deste modo podemos comunicar de maneira mais ágil.
            \item \textbf{Análise de Documentação} como \textbf{artigos de jornais} sobre casos resolvidos pela consultoria, para a equipa estar mais contextualizada acerca da mesma.
            \textbf{Relatórios de casos resolvidos} são recursos escassos infelizmente, pelo que também recorremos à \textbf{leitura de entrevistas dos detetives} para uma melhor compreensão das estratégias e métodos que estes têm por hábito utilizar na sua resolução dos seus casos. Analisámos igualmente a \textbf{documentação} da própria consultoria, mais especificamente relatórios relacionados com os ganhos e custos ao longo dos anos da empresa.
            \end{itemize}
            
            \par Durante a reunião foi levantada a possibilidade de acompanhar uma investigação de perto, infelizmente foi negada por questões de segurança e de confidencialidade das informações dos participantes da investigação. 

            
    \section{Organização dos Requisitos Levantados}
        \textcolor{red}{
            >> Explicação do processo de levantamento realizado e exposição da organização de requisitos adotada. Enumeração e organização dos requisitos levantados, de acordo com a sua categorização (descrição, manipulação e controlo). Caracterização dos requisitos tendo em consideração os seguintes elementos: tipo, número, data, descrição do requisito, fonte de informação e analista.
        }
    \section{Análise e Validação Geral dos Requisitos}
        \textcolor{red}{
            >> Validação dos requisitos apresentados com os diversos intervenientes do processo. Relatar eventuais anomalias e ações corretivas desenvolvidas.
        }

%==========================================================================
% END LEVANTAMENTO E ANÁLISE DE REQUISITOS
%==========================================================================

%==========================================================================
% BEGIN MODELAÇÃO CONCEPTUAL
%==========================================================================

\chapter{Modelação Conceptual}
    \section{Apresentação da Abordagem de Modelação Realizada}
        \textcolor{red}{
            <<Apresentação e explicação do processo de modelação concetual adotado, com referência à notação e ferramenta adotadas.>>
        }
    \section{Identificação e Caracterização das Entidades}
        \textcolor{red}{
            <<Explicação do processo de identificação das entidades do sistema. Enumeração e fundamentação de cada uma das entidades identificadas, tendo em conta a caracterização sugerida em [Connoly e Begg, 2015]. Indicar os requisitos que originaram cada uma das entidades.>>
        }
    \section{Identificação e Caracterização dos Relacionamentos}
        \textcolor{red}{
            <<Explicação do processo de identificação dos relacionamentos entre as entidades definidas. Enumeração e fundamentação de cada um dos relacionamentos identificados, tendo em conta a caracterização sugerida em [Connolly e Begg, 2015]. Indicar os requisitos que originaram cada um dos relacionamentos enumerados.>>
        }
    \section{Identificação e Caracterização dos Atributos das Entidades e dos Relacionamentos}
        \textcolor{red}{
            <<Explicação do processo de identificação dos vários atributos das entidades e dos relacionamentos. Enumeração e fundamentação de cada atributo, organizando-os por entidades e relacionamentos identificados, tendo em conta a caracterização sugerida em [Connolly e Begg, 2015]. Indicar os requisitos que originaram cada um dos atributos enumerados.>>
        }
    \section{Apresentação e Explicação do Diagrama ER Produzido}
        \textcolor{red}{
            <<Apresentação do diagrama concetual produzido. Explicação do seu processo de construção, tendo em conta as entidades e os relacionamentos identificados anteriormente, bem com as suas caracterizações em termos de atributos.>>
        }

%==========================================================================
% END MODELAÇÃO CONCEPTUAL
%==========================================================================

%==========================================================================
% BEGIN MODELAÇÃO LÓGICA
%==========================================================================

\chapter{Modelação Lógica}
    \section{Construção e Validação do Modelo de Dados Lógico}
        \textcolor{red}{
            <<Apresentação e explicação do processo de modelação lógica adotado, com referência à ferramenta adotada.>>
        }
    \section{Apresentação e Explicação do Modelo Lógico Produzido}
        \textcolor{red}{
            <<Apresentação do processo de conversão realizado, expondo e justificando a origem de cada uma das tabelas que foram criadas. Apresentação do modelo lógico produzido.>>
        }
    \section{Normalização de Dados}
        \textcolor{red}{
            <<Indicação se o modelo está ou não normalizado. Explicar.>>
        }
    \section{Validação do Modelo com Interrogações do Utilizador}
        \textcolor{red}{
            <<Apresentação de 4-6 expressões em Álgebra Relacional que representem queries previamente enunciadas no conjunto de requisitos de manipulação estabelecidos anteriormente.>>
        }

%==========================================================================
% END MODELAÇÃO LÓGICA
%==========================================================================

%==========================================================================
% BEGIN CONCLUSÕES DE TRABALHO FUTURO
%==========================================================================

\chapter{Conclusões e Trabalho Futuro}
    \textcolor{red}{
        <<Elaborar uma apreciação crítica sobre o trabalho realizado, apontando os seus pontos fortes e fracos. Adicionalmente, caso se aplique, enunciar eventuais tarefas a realizar futuramente ou novas opções para estender o trabalho realizado.>> \\
        <<Resumo breve do trabalho realizado e das ações desenvolvidas. Enumeração e análise de aspetos positivos e negativos identificados durante o processo de desenvolvimento do sistema de bases de dados. Exposição das próximas linhas de desenvolvimneto do projeto.>>
    }

%==========================================================================
% END CONCLUSÕES DE TRABALHO FUTURO
%==========================================================================

%==========================================================================
% BEGIN BIBLIOGRAFIA
%==========================================================================

%% Changes biblibography name
%% Portuguese babel default : "Bibliografia"
%% Personally I prefer "Referências"
% \renewcommand\bibname{Referências}

%% https://www.overleaf.com/learn/latex/bibliography_management_with_bibtex
\begin{thebibliography}{9}
\textcolor{red}{
    <<Enumeração dos diversos recursos bibliográficos utilizados. Utilizar o sistema de referenciação APA (https://guias.sdum.uminho.pt/apa).>>
}
\end{thebibliography}

%==========================================================================
% END BIBLIOGRAFIA
%==========================================================================

%==========================================================================
% BEGIN LISTA DE SIGLAS E ACRÓNIMOS
%==========================================================================

%% Portuguese babel does not translate this environment
\renewcommand{\nomname}{Lista de Siglas e Acrónimos}

%% Text that can be shown before acronyms list
% \renewcommand{\nompreamble}{
%     \textcolor{red}{
%         <<Apresentar uma lista com todas as siglas e acrónimos utilizados durante a realização do trabalho. O formato base para esta lista deverá ser da forma como abaixo se apresenta.>>
%     }
% }

%% acronyms
\nomenclature[01]{\textbf{CDC}}{Consultoria de Detetives Christie}
\nomenclature[02]{\textbf{BD}}{Base de Dados}
\nomenclature[03]{\textbf{SBD}}{Sistema de Base de Dados}
\nomenclature[04]{\textbf{ER}}{Entidade-Relacionamento}

%% Show acronyms
\printnomenclature

%==========================================================================
% END LISTA DE SIGLAS E ACRÓNIMOS
%==========================================================================

%==========================================================================
% BEGIN ANEXOS
%==========================================================================

%% Why \addchap, instead of \chapter?
%% \addchap has no numbering but appears in table of contents.
\addchap{Anexos}

    \textcolor{red}{
        <<Os anexos deverão ser utilizados para a inclusão de informação adicional necessária para uma melhor compreensão do relatório o para complementar tópicos, secções ou assuntos abordados. Os anexos criados deverão ser numerados e possuir uma designação. Estes dados permitirão complementar o Índice geral do relatório relativamente à enumeração e apresentação dos diversos anexos.>>
    }

    %% section version of \addchap
    \addsec{Anexo 1}


%==========================================================================
% END ANEXOS
%==========================================================================
\end{document}
